\chapter{Why this matters, and what I take from it}

I said at the start that I did not want to write an attack, and I meant it. Rerunning someone's analysis, finding that it does not hold, and explaining why, is ordinary science. It is how the consensus in Figure~\ref{fig:consensus_by_field} was built in the first place: not by people agreeing with each other, but by more than a century of people trying to prove each other wrong and failing. The radiative physics of CO₂ was measured in laboratories in the 1850s. The records were built, found to contain errors, and corrected. The models were tested against ice cores, ocean heat and satellite measurements of the energy leaving the planet. Agreement of 97\% to 99\% among the people who do that work is what convergence looks like.

What this case teaches, I think, is that the climate is a system, and you need to know its parts before you can read its signals. A competent statistician, coming to it from outside, can easily miss the aerosols, the oceans, the station corrections, and the difference between yearly and accumulated emissions, and end up with a clean-looking chart that points the wrong way. The fix is not to tell people from other fields to stay out. It is to read the field's literature first, to talk to the people who built the data, and to be most careful precisely when a result overturns what the field has found.

And there is a reason this matters beyond the seminar room. Doubt about climate change has never been evenly spread, and neither has money. The fields with the most doubt in Figure~\ref{fig:consensus_by_field} are the fields closest to extraction. The companies on Heede's list spent decades funding public confusion about the science after their own researchers had confirmed the warming. At a conference whose keynotes were about building more data centres to power artificial intelligence, a talk saying that CO₂ is not the main driver of warming will be welcome to some in the room, whatever its author intended. The electricity those data centres need will come from somewhere, and whether it comes from coal and gas or from wind and sun is exactly the kind of decision that a clear understanding of the science should inform.

So read the thermometer carefully, and then read the ledger. The thermometer shows that the world has warmed faster in the past four decades than at any time in the instrumental record, and that the warming tracks the CO₂ we have accumulated. The ledger shows that this is not a story about ``humanity'' in general, but about a particular way of producing things, one that made burning carbon profitable and kept it profitable long after its cost was known. Systems built by people can be changed by people, and that work starts with getting the science right and explaining it patiently to anyone who is willing to look.

\begin{center}\rule{0.5\linewidth}{0.5pt}\end{center}
